\documentclass[a4paper,12pt]{ltjsarticle}
\usepackage[top=25mm,bottom=25mm,left=25mm,right=25mm]{geometry}
\usepackage{booktabs}
\usepackage{array}
\usepackage{longtable}
\usepackage{enumitem}
\usepackage{hyperref}
\usepackage{titlesec}
\usepackage{xcolor}
\usepackage{graphicx}

\setlist[itemize]{leftmargin=2em}
\setlist[itemize,2]{leftmargin=2em}
\setlist[enumerate]{leftmargin=2em}
\setlist[enumerate,2]{leftmargin=2em}

\newcolumntype{L}[1]{>{\raggedright\arraybackslash}p{#1}}
\newcolumntype{C}[1]{>{\centering\arraybackslash}p{#1}}

% セクション見出しのスタイル
\titleformat{\section}{\large\bfseries}{第\thesection 節}{1em}{}[\titlerule]
\titleformat{\subsection}{\normalsize\bfseries}{\thesubsection}{1em}{}

\begin{document}

% ヘッダ情報
\begin{center}
  {\LARGE\bfseries 作業ログ}
\end{center}

\vspace{4mm}

\begin{tabular}{ll}
  \textbf{報告者} & 酒井 睦基 \\
  \textbf{報告日} & \today \\
  \textbf{プロジェクト名} & Arduino オーケストラ：きらきら星 輪唱演奏システム \\
\end{tabular}

\vspace{6mm}

%-------------------------------------------------------------------
\section{進捗サマリー（要約）}

\begin{itemize}
  \item \textbf{全体進捗率}：30 [\%]
  \item \textbf{マイルストーン達成状況}：未達成
  \item \textbf{主要成果}：
    \begin{itemize}
      \item 共通ライブラリの整備及び内容の共有が完了
    \end{itemize}
  \item \textbf{未達成要項}：
    \begin{itemize}
      \item マスタ1台とスレーブ4台の同期の確立を目的としていたが中間テスト対策のため未達
      \item ベビーメリー用の材料調達が未達
    \end{itemize}
  \item \textbf{重大課題（影響度：高）}：ベビーメリーとBPM用の材料を3年が準備する必要あり
  \item \textbf{重大課題（影響度：中）}：全5台の同期の確立を実装する必要あり（ガントチャート上は第7〜9週で行うため進捗としては通常）
  \item \textbf{重大課題（影響度：低）}：各楽器音の細かな修正の余地あり
\end{itemize}

%-------------------------------------------------------------------
\section{作業ログ（詳細トラッキング）}

\begin{longtable}{C{2cm}C{1.6cm}C{1.6cm}C{2.2cm}C{1.2cm}L{2.4cm}L{3cm}}
  \toprule
  \textbf{作業項目} & \textbf{開始日} & \textbf{予定完了日} & \textbf{状態} & \textbf{進捗率} & \textbf{依存関係・ブロッカー} & \textbf{コメント} \\
  \midrule
  \endfirsthead
  \toprule
  \textbf{作業項目} & \textbf{開始日} & \textbf{予定完了日} & \textbf{状態} & \textbf{進捗率} & \textbf{依存関係・ブロッカー} & \textbf{コメント} \\
  \midrule
  \endhead
  \midrule
  \endfoot
  \bottomrule
  \endlastfoot
  共通ライブラリの作成 & 5/20 & 6/2 & 完了 & 100\% & 2年生3人が協力してライブラリをまとめ，内容も3人の間で共有されていることを確認 & - \\
\end{longtable}

%-------------------------------------------------------------------
\section{技術成熟度レベルTRLの推移} \label{trl}
%-------------------------------------------------------------------
FBSの各要素ごとに現在どのレベルにいるかを以下に示す．

\begin{longtable}{L{2.5cm}L{4cm}C{2cm}C{2cm}L{3.5cm}}
  \toprule
  \textbf{機能} & \textbf{説明} & \textbf{前回TRL} & \textbf{今回TRL} & \textbf{備考} \\
  \midrule
  \endfirsthead
  \toprule
  \textbf{機能} & \textbf{説明} & \textbf{前回TRL} & \textbf{今回TRL} & \textbf{備考} \\
  \midrule
  \endhead
  \midrule
  \endfoot
  \bottomrule
  \endlastfoot
  I2C同期（F1） & マスタ→スレーブ同期通信 & 1 & 2 & マスタ1台，スレーブ2台の接続を確認できた \\
  音色合成（F4） & Processing加算合成 & 1 & 2 & 4種類の楽器音が確認できた \\
  PLL補正（F1.3） & 位相同期制御 & 1 & 1 & 設計中 \\
  エンタメ機能（F5） & ロードセル・サーボ制御 & 1 & 1 & 設計中 \\
  可変テンポ（F2） & BPM制御 & 1 & 1 & 設計中 \\
  共通ライブラリ（F3.1） & PROGMEM＋I2Cフレーム & 1 & 2 & 共通ライブラリを確認できた \\
\end{longtable}

\noindent
{\small
\begin{tabular}{ll}
  \textbf{TRL} & \textbf{定義} \\
  1 & 概念レベル：アイディアや原理が確立されており，説明できるレベル \\
  2 & 基本動作レベル：単体で動作確認が可能なレベル \\
  3 & 結合動作レベル：機能を結合し，実際の使用環境に近い環境で動作確認が可能なレベル \\
  4 & 実運用レベル：実際の使用環境で安定的な稼働が確認できるレベル \\
\end{tabular}
}

%-------------------------------------------------------------------
\section{技術性能指標TPMの推移} \label{tpm}
%-------------------------------------------------------------------
MOPの各要素毎に現在の値の程度を以下に示す．

\begin{longtable}{L{2cm}L{3cm}C{2cm}C{2cm}C{2cm}L{2.5cm}}
  \toprule
  \textbf{関連MOP} & \textbf{指標} & \textbf{目標値} & \textbf{前回値} & \textbf{今回値} & \textbf{備考} \\
  \midrule
  \endfirsthead
  \toprule
  \textbf{関連MOP} & \textbf{指標} & \textbf{目標値} & \textbf{前回値} & \textbf{今回値} & \textbf{備考} \\
  \midrule
  \endhead
  \midrule
  \endfoot
  \bottomrule
  \endlastfoot
  MOP1 & 楽器間同期誤差 & $\leq 30$\,ms & --- & --- & 未計測 \\
  MOP2 & I2C SYNCパケットロス率 & 0回 & --- & --- & 未計測 \\
  MOP3 & 楽曲識別（中間確認） & チーム内識別 & --- & --- & 未実施 \\
  MOP4 & 楽器識別（中間確認） & チーム内識別 & --- & --- & 未実施 \\
  MOP5 & リラクゼーション感 & 心地よく聞こえる & --- & --- & 未実施 \\
  MOP6 & BPM変更反映速度 & 1小節以内 & --- & --- & 未計測 \\
  MOP7 & PLL収束速度 & 変化量$\div$10小節以内 & --- & --- & 未計測 \\
\end{longtable}

\subsection{現在のガントチャート}
現在のガントチャートの状態を図\ref{fig:gc}に示す．
灰色の箇所は実装が完了を示す．

\begin{figure}[htbp]
  \centering
  \includegraphics[width=\linewidth]{fig/new_gant.png}
  \caption{現在のガントチャート}
  \label{fig:gc}
\end{figure}

\subsection{日毎の作業時間と作業内容}
定期的にGitHubを確認しPRがきていないかどうかを確認する．
PRがきていた場合はファイルの内容をレビューし，問題がない場合はmergeする．
基本的な作業時間はGitHubの確認とコードレビューを含め1時間程度である．
また，前回のMTGでZulipのコメント数に関する指摘を受けたため，SlackのコメントをZulipに移行した．
移行にかかった時間は30分程度である．

%-------------------------------------------------------------------
\section{課題管理}
\begin{itemize}
  \item \textbf{課題}：筐体設計
    \begin{itemize}
      \item \textbf{発生原因}：材料の確保と設計が必要
      \item \textbf{影響度}：高
      \item \textbf{優先度}：高
      \item \textbf{対策}：3年生が主体となり6/10のMTGで筐体設計に必要な材料確保の日程と設計図の作成を行う
      \item \textbf{期限}：6/9
      \item \textbf{状態}：未対応
    \end{itemize}
\end{itemize}

% \noindent
% {\small ※影響度＝プロジェクトへのインパクト \quad ※優先度＝対応の緊急性}

%-------------------------------------------------------------------
\section{AI利用状況と検証（再現性重視）}
TRLとTPMのテンプレート作成のみに利用したため，検証方法，ファクトチェック結果の記載を省く．

\subsection{利用内容}
\begin{itemize}
  \item \textbf{利用目的}：TRLとTPMセクションのテンプレート作成
  \item \textbf{入力（プロンプト要旨）}：このtexファイルにTRLとTPMのセクションを追加して
  \item \textbf{出力概要}：セクション\ref{trl}，\ref{tpm}が作成された
\end{itemize}

\subsection{検証方法}

\begin{itemize}
  \item \textbf{検証手法}：
    \begin{itemize}
      \item 実験／比較／理論確認など
    \end{itemize}
  \item \textbf{参照資料}：
    \begin{itemize}
      \item 論文・書籍・公式ドキュメント等
    \end{itemize}
  \item \textbf{検証手順}：
    \begin{enumerate}
      \item 手順1
      \item 手順2
    \end{enumerate}
\end{itemize}

\subsection{ファクトチェック結果}
\begin{itemize}
  \item \textbf{一致点}：
    \begin{itemize}
      \item 既存知識と整合している部分
    \end{itemize}
  \item \textbf{不一致・疑義}：
    \begin{itemize}
      \item 差異の内容
    \end{itemize}
  \item \textbf{最終判断}：採用／保留／棄却
  \item \textbf{根拠}：
    \begin{itemize}
      \item 資料・理由
    \end{itemize}
\end{itemize}

%-------------------------------------------------------------------
\section{次回計画}

\begin{itemize}
  \item \textbf{3年生}：
    \begin{itemize}
      \item WBS 1.3：ベビーメリー風装置の設計・製作（全員）
    \end{itemize}
  \item \textbf{2年生}：
    \begin{itemize}
      \item WBS 2.1.6：Serial通信・制御ロジック実装（成毛）
      \item WBS 2.2，2.2.1：BPM変化による可変テンポ機能を実装（成毛）
      \item WBS 2.2.3，2.2.4：マスタ側同期制御実装（佐藤・山口）
      \item WBS 2.3.1：スレーブ側で音符データの受け取りを実装（佐藤・山口）
    \end{itemize}
  \item \textbf{期限}：6/16（第8週末）
\end{itemize}

%-------------------------------------------------------------------
\section{その他}

\begin{itemize}
  \item 特になし
\end{itemize}

%-------------------------------------------------------------------
\section{付録}

\begin{itemize}
  \item \textbf{添付資料}：
    \begin{itemize}
      \item GitHub（ガントチャートの各タスク毎にGitHubでissueを作成した，各タスクの完了報告はissueにて確認している）

        \begin{figure}[htbp]
            \centering
            \includegraphics[width=\linewidth]{fig/issue_1.png}
            \caption{GitHubのissueの例1}
            \label{fig:is1}
        \end{figure}

        \begin{figure}[htbp]
            \centering
            \includegraphics[width=\linewidth]{fig/issue_2.png}
            \caption{GitHubのissueの例2}
            \label{fig:is2}
        \end{figure}
        
    \end{itemize}

  \item \textbf{添付範囲}：全部
\end{itemize}

\end{document}